\section{Heterotic line-bundle models}
\label{sec:setup}

\subsection{Geometry}

Let $A=\PP^{n_1}\times\cdots\times\PP^{n_m}$ and let $X\subset A$ be a smooth
complete intersection Calabi--Yau threefold~\cite{CDLS1988,GHL1989,Hubsch1992}. We
assume $X$ is \emph{favourable}: the restrictions $J_1,\dots,J_h$ of the
hyperplane classes of the factors form an integral basis of $H^{1,1}(X)$, so
$h=m$ and every line bundle is $\Oc_X(k)\coloneqq\Oc_A(k)|_X$ for a unique
$k\in\Z^h$. The topological data are the triple intersection numbers and the
second Chern class,
\begin{equation}
d_{ijl}=\int_X J_iJ_jJ_l\in\Z,\qquad c_{2,i}=\int_X c_2(TX)\,J_i ,
\end{equation}
both computed exactly from the ambient Chow ring and the tangent--normal
sequence. We write
\begin{equation}\label{eq:notation}
D(x,y,z)=d_{ijl}x^iy^jz^l,\qquad
q_r(k)=d_{rij}k^ik^j,\qquad
M(k)_{jl}=d_{ijl}k^i,
\end{equation}
so that $M:\Z^h\to\Sym_h(\Z)$ is a group homomorphism. By
Hirzebruch--Riemann--Roch the index of a line bundle is
\begin{equation}\label{eq:hrr}
\ind\Oc_X(k)=\sum_q(-1)^q h^q(X,\Oc_X(k))=\tfrac16D(k,k,k)+\tfrac1{12}c_{2,i}k^i .
\end{equation}

\begin{table}[t]
\centering
\small
\begin{tabular}{@{}lcc@{}}
\toprule
 & $\Xa$ (tetraquadric) & $\Xb$ \\
\midrule
ambient space & $(\PP^1)^4$ & $(\PP^1)^5$ \\
defining equations & one of multidegree $(2,2,2,2)$ & two of multidegree $(1,1,1,1,1)$ \\
$(h^{1,1},h^{2,1})$ & $(4,68)$ & $(5,45)$ \\
$d_{ijl}$ & $2$ for $i,j,l$ distinct, else $0$ & $2$ for $i,j,l$ distinct, else $0$ \\
$c_{2,i}$ & $24$ for all $i$ & $24$ for all $i$ \\
index & $\ind\Oc_X(k)=2\bigl(e_3(k)+e_1(k)\bigr)$ & $\ind\Oc_X(k)=2\bigl(e_3(k)+e_1(k)\bigr)$ \\
$d_X=\gcd\ind$ & 2 & 2 \\
anomaly constant $C_X$ (\cref{eq:CX}) & 48 & 40 \\
free actions used & $\Z_2$, $\Z_2\times\Z_2$ & $\Z_2$ \\
\bottomrule
\end{tabular}
\caption{The two benchmark geometries. $e_j$ denotes the $j$-th elementary
symmetric polynomial in the components of $k$. Topology was recomputed from the
configuration matrices and agrees with the published values for all 68
favourable CICYs ($c_2$: 68/68, Euler number: 68/68).}
\label{tab:geometries}
\end{table}

\subsection{Line-bundle sums and the upstairs conditions}

Let $V=\bigoplus_{a=1}^5 L_a$, $L_a=\Oc_X(k_a)$. With $V$ embedded in one $E_8$
factor, its structure group $S(U(1)^5)\subset SU(5)$ leaves an $SU(5)$ GUT times
$U(1)^4$ in four dimensions. The upstairs conditions
of~\cite{ACGLP2013} on the charge matrix $K=(k_1,\dots,k_5)$ are:

\begin{enumerate}[label=(C\arabic*),leftmargin=*]
\item\label{c:c1} $c_1(V)=\sum_a k_a=0$.
\item\label{c:structure} No nonempty proper subset $S\subsetneq\{1,\dots,5\}$ has
  $\sum_{a\in S}k_a=0$, so the structure group is not smaller than $S(U(1)^5)$.
\item\label{c:index} $\ind V=\ind\wedge^2V=-3|\Gam|$, giving three chiral
  families after the quotient by a freely acting $\Gam$.
\item\label{c:anomaly} $c_2(TX)-c_2(V)$ is effective (trivial hidden bundle).
  Since the $J_r$ are nef, a necessary condition is
  $c_{2,r}-c_2(V)_r\ge0$ for all $r$. With \ref{c:c1},
  $c_2(V)_r=-\tfrac12\sum_a q_r(k_a)$.
\item\label{c:slope} $V$ is poly-stable: there is a Kähler class $J=t^iJ_i$ at
  which all five slopes $\mu_{k_a}(t)=t^\top M(k_a)t$ vanish simultaneously.
\item\label{c:bogomolov} Bogomolov: $c_2(V)\cdot J\ge 0$ at that class.
\item\label{c:spectrum} For the target \st{GUT\_FULL}: $h^\bullet(X,V)=(0,3|\Gam|,0,0)$
  and $h^\bullet(X,\wedge^2V)=(0,3|\Gam|+n_h,n_h,0)$ with $n_h>0$ (no
  anti-families; at least one Higgs pair).
\end{enumerate}

Throughout, the Kähler parameter is restricted to the positive ambient chamber
$T_+=\{t\in\R^h:t^i>0\}$; for both benchmarks this is the declared Kähler cone,
an input assumption carried in every certificate.

\subsection{Descent and the downstairs spectrum}

Let $\Gam$ be a freely acting abelian group, $\hat X=X/\Gam$. An equivariant
structure on each $L_a$ lets $V$ descend; a Wilson line, i.e.\ a homomorphism
$\Gam\to U(1)_Y\subset SU(5)$, breaks $SU(5)\to SU(3)\times SU(2)\times U(1)_Y$.
Upstairs, $H^1(X,V)$ carries the $\ten$s, $H^1(X,\wedge^2V)$ the $\bfive$s,
$H^2(X,\wedge^2V)\cong H^1(X,\wedge^2V^*)^*$ the $\five$s, and
$H^1(X,L_a\otimes L_b^*)$ the singlets $S_{ab}$. Downstairs fields are the
$\Gam$-invariant parts of these cohomologies tensored with the Wilson-line
representation. The target \st{DOWNSTAIRS\_MSSM} requires three chiral families
of $Q,u^c,e^c,d^c,L$, no charged chiral exotics, and at least one Higgs doublet
pair; the distinct stricter target of an exact non-supersymmetric Standard Model
is not addressed.

\subsection{The search problem}

For a bound $K\in\Z_{>0}$ let
\begin{equation}
\mathcal D_K=\bigl\{K\in\Z^{h\times5}\;:\;|k_a^i|\le K,\ \textstyle\sum_ak_a=0\bigr\}/S_5 ,
\end{equation}
the charge matrices with entries in the box, up to permutation of the summands.
The \emph{bounded search problem} is to determine the subset of $\mathcal D_K$
satisfying \ref{c:c1}--\ref{c:spectrum}. Its nominal size is enormous: already at
$K=7$ on $\Xa$ there are $15^{16}\approx6.6\times10^{18}$ ordered choices of the
first four columns. The \emph{unbounded problem}---to decide the full solution
set in $\Z^{h\times5}$---is the real question, and \cref{sec:slices} is our
partial answer to it.
